% !TEX root = ../neurips_2026.tex

\section{Introduction}
\label{sec:intro}

Knowledge distillation (KD) compresses a large \emph{teacher} language model into a smaller \emph{student} by training the student to match the teacher's output distribution~\citep{hinton2015distilling}.
For autoregressive LLMs, this amounts to minimizing the per-token KL divergence between teacher and student next-token distributions across a training corpus.
Recent advances have refined the divergence measure~\citep{gu2024minillm, agarwal2024gkd, ko2024distillm}, introduced adaptive loss functions~\citep{wu2025akl}, and explored difficulty-aware curricula~\citep{dakd2025}, steadily closing the quality gap between teacher and student.

Yet all these methods share a fundamental limitation: they constrain \emph{what} the student predicts but not \emph{why}.
Consider a reading comprehension task where the teacher correctly answers ``Paris'' by attending to the sentence ``The capital of France is Paris.''
A student trained with standard KD may produce the same answer by latching onto the keyword ``France'' in the question---a shallow shortcut that yields identical output KL yet fails on rephrased or adversarial inputs where the shortcut breaks.
The crux is that output-level matching provides no guarantee about which input tokens drive the prediction.

\paragraph{The Jacobian gap.}
We formalize this observation through function approximation theory.
Let $f_T$ and $f_S$ denote the teacher and student mappings from input embeddings to output log-probabilities.
Standard KD minimizes $D_{\kl}(f_T(x) \| f_S(x))$ at each training point $x$---a \emph{zero-order} (pointwise) constraint.
A first-order Taylor expansion at a perturbed input $x + \delta$ ($\|\delta\| \leq \epsilon$) reveals:
\begin{equation}
  D_{\kl}\bigl(f_T(x{+}\delta) \,\|\, f_S(x{+}\delta)\bigr)
  \;\leq\;
  \underbrace{D_{\kl}\bigl(f_T(x) \,\|\, f_S(x)\bigr)}_{\text{zero-order: standard KD}}
  \;+\; \epsilon \cdot
  \underbrace{\|J_T(x) - J_S(x)\|_F}_{\text{first-order: Jacobian gap}}
  \;+\; O(\epsilon^2),
  \label{eq:taylor-intro}
\end{equation}
where $J_T, J_S$ are the input--output Jacobians.
Standard KD drives the first term to zero but leaves the Jacobian gap entirely unconstrained.
In approximation-theoretic terms, KD achieves $L^2$ convergence but not \emph{Sobolev} $W^{1,2}$ convergence---the latter being necessary to control the student's behavior in the input neighborhood~\citep{czarnecki2017sobolev}.

\paragraph{Why not just match the Jacobian?}
For an LLM with vocabulary size $V$, sequence length $L$, and hidden dimension $d$, the full Jacobian $J \in \R^{V \times (L \cdot d)}$ has on the order of $10^{11}$ entries.
Explicit computation, storage, or alignment is intractable.
Prior work on Jacobian matching in KD has been limited to small encoder-only models~\citep{srinivas2018knowledge, wang2022gkd} or restricted to classification tasks where the output is a single vector~\citep{wu-etal-2023-ad}.
No existing method achieves Jacobian-level matching for decoder-only LLM generation.

\paragraph{Key insight: noise as an implicit Jacobian probe.}
We build on a classical result by \citet{srinivas2018knowledge}: for MSE loss on classification, the expected loss under Gaussian input noise recovers the Jacobian Frobenius norm.
We extend this principle to the autoregressive, token-level KL setting of LLM distillation.
Specifically, we show (Theorem~\ref{thm:noise-identity}) that
\begin{equation}
  \E_{\xi \sim \mathcal{N}(0, \sigma^2 I)}\bigl[D_{\kl}\bigl(f_T(x{+}\xi) \,\|\, f_S(x{+}\xi)\bigr)\bigr]
  \;=\;
  D_{\kl}\bigl(f_T(x) \,\|\, f_S(x)\bigr)
  \;+\; \frac{\sigma^2}{2} \sum_{t \in \cR} \cF_t(x)
  \;+\; O(\sigma^4),
  \label{eq:noise-identity-intro}
\end{equation}
where $\cF_t(x) = \sum_v p_T^v \|\nabla_x \log p_T^v - \nabla_x \log p_S^v\|^2$ is a Fisher-weighted score matching term that captures the per-position Jacobian difference, and $\cR$ denotes response token positions.
By simply evaluating KL on noise-perturbed embeddings, we \emph{implicitly} match the full Jacobian---without computing it, without second-order gradients, and without requiring teacher and student to share the same hidden dimension.

\paragraph{Saliency-Guided Knowledge Distillation (\sagd{}).}
Building on this insight, we propose \sagd{}, which augments standard KD with two complementary mechanisms:

\begin{enumerate}
\item \textbf{Noise-injected KL (implicit Jacobian matching).}
  At every $N$-th training step, we add position-adaptive Gaussian noise to input embeddings and compute KL between teacher and student on the perturbed inputs.
  The noise scale at each input position is proportional to the local saliency disagreement between teacher and student, concentrating the implicit Jacobian probe where it matters most.
  This upgrades the distillation guarantee from $L^2$ to Sobolev $W^{1,2}$.

\item \textbf{Saliency-guided sample reweighting (DRO).}
  We compute \emph{input saliency}---the $\ell_2$-norm of the embedding gradient of the response log-probability---for both teacher and student, and use their Jensen--Shannon divergence as a per-sample difficulty measure.
  Samples where teacher and student attend to different tokens receive higher training weight.
  We show this corresponds to entropy-regularized distributionally robust optimization that minimizes the worst-case neighborhood error across the training set (Theorem~\ref{thm:dro}).
\end{enumerate}

Neither mechanism alone is sufficient.
Noise KL tightens the Jacobian bound uniformly across all samples; reweighting allocates more optimization budget to samples with the loosest bounds.
Together, they ensure both that the student's input sensitivity matches the teacher's (first-order) and that this matching is prioritized where it is most needed (robust allocation).

\paragraph{Evidence concentration: a diagnostic beyond outputs.}
Standard evaluation of distilled models relies entirely on output-level metrics (ROUGE-L, Exact Match, perplexity).
We introduce \emph{evidence concentration}---the fraction of the student's input saliency that falls on the ground-truth answer span---as a diagnostic for whether the student \emph{reasons} like the teacher.
On SQuAD~2.0, we find that standard KD students dramatically over-concentrate saliency on answer tokens (EC\,$=$\,0.169 vs.\ teacher's 0.055), indicating shortcut learning.
\sagd{} reduces this to 0.083, bringing the student's reasoning pattern significantly closer to the teacher's holistic evidence utilization.

\paragraph{Contributions.}
\begin{enumerate}
\item \textbf{Theory.}
  We derive the noise--Jacobian identity for autoregressive token-level KL (extending \citealt{srinivas2018knowledge} beyond MSE/classification), establish a neighborhood error bound connecting noise KL to Sobolev $W^{1,2}$ convergence (Theorem~\ref{thm:neighborhood}), and formalize saliency-guided reweighting as entropy-regularized DRO on the Jacobian gap (Theorem~\ref{thm:dro}).

\item \textbf{Method.}
  We propose \sagd{}, combining noise-injected KL for implicit Jacobian matching with saliency-guided sample reweighting.
  The method requires no architectural constraints, no second-order gradients, and adds ${\sim}40\%$ overhead relative to standard KD (with saliency computed every $N{=}5$ steps).

\item \textbf{Experiments.}
  On SQuAD~2.0 and Dolly-15K with Qwen3-8B $\to$ 0.6B/1.7B, \sagd{} improves Exact Match, Token F1, and ROUGE-L over standard KD and six recent baselines.
  Evidence concentration analysis confirms that improvements stem from preserving the teacher's internal reasoning, not merely output mimicry.
  Ablations isolate the contributions of noise KL and reweighting, and hyperparameter sweeps validate robustness.
\end{enumerate}
